\cleardoublepage
\chapter{Podsumowanie}

W~ramach niniejszej pracy:

\begin{itemize}
\item zebrano i~usystematyzowano wiedzę z~zakresu architektury mikroserwisowej, mechanizmów komunikacji synchronicznej i~asynchronicznej oraz wzorców projektowych dla systemów rozproszonych w~ekosystemie .NET,
\item dokonano szczegółowej analizy protokołu gRPC-Web i~jego wariantów wdrożeniowych (proxy Envoy, bezpośrednia obsługa w~oprogramowaniu pośredniczącym ASP.NET Core),
\item zaprojektowano i~zaimplementowano autorski system testowy składający się z~dwóch mikroserwisów .NET~10 (ProductService, OrderService), infrastruktury wspierającej (PostgreSQL, Redis, Envoy Proxy) oraz warstwy obserwowalności (Prometheus, Grafana),
\item zaimplementowano wzorzec Cache-Aside z~wykorzystaniem Redis, a~także skonfigurowano serwer Kestrel do~równoległej obsługi trzech protokołów,
\item opracowano metodykę pomiarową opartą na~rzeczywistej przeglądarce internetowej, w~której czas mierzony jest wewnątrz strony, od~interakcji użytkownika do~naniesienia odpowiedzi na~drzewo dokumentu, co~obejmuje pełną ścieżkę odczuwaną przez użytkownika aplikacji przeglądarkowej,
\item przeprowadzono 39~konfiguracji pomiarowych w~pięciu niezależnych powtórzeniach z~rotowaną kolejnością protokołów, co~dało 585~wykonanych testów i~171~tysięcy pomiarów, badając wpływ protokołu komunikacji, rozmiaru odpowiedzi, stanu pamięci podręcznej oraz liczby równoległych klientów na~czas odpowiedzi systemu,
\item poddano wyniki analizie statystycznej, wyznaczając przedziały ufności dla różnic sparowanych w~obrębie przebiegu.
\end{itemize}

Główne wnioski z~przeprowadzonych badań eksperymentalnych:

\begin{enumerate}
\item Warianty gRPC-Web skracają czas odpowiedzi o~26--43\% względem REST dla dużych, nieskompresowanych odpowiedzi. Efekt pojawiał się po~przekroczeniu zaobserwowanego w~badanym środowisku progu około 64~KiB. Poniżej tego progu różnice wynosiły kilka procent.

\item W~analizie objętościowej gzip zmniejszył dokument JSON dwunasto- do~trzynastokrotnie, a~wiadomość binarną około dziesięciokrotnie. Przewaga objętościowa formatu binarnego spadła z~33\% do~15--18\%, a~skompresowane ładunki nie przekroczyły progu 64~KiB. Czasu odpowiedzi z~włączoną kompresją nie mierzono.

\item Zależność przewagi od~rozmiaru odpowiedzi miała charakter skokowy. Model oparty na~efektywnym progu transmisji około 64~KiB odtworzył zmierzone przyrosty czasu w~dziesięciu konfiguracjach z~dokładnością do~0,29~RTT.

\item W~pięciu konfiguracjach dekompozycji różnica REST--Envoy stanowiła około 85--98\% całkowitej różnicy REST--Direct; dla 500~produktów i~dziesięciu klientów udział wynosił około 95\%. Wyniki wspierają znaczenie redukcji wolumenu danych. Porównanie kompletnych ścieżek komunikacji daje przybliżony podział kosztów serializacji i~transportu.

\item Serializacja Protocol Buffers redukuje wolumen przesyłanych danych o~32--43\%. Wielkość ta~nie zależy od~warunków sieciowych, obciążenia ani od~zastosowanego narzędzia pomiarowego. Jedynym czynnikiem, który ją~zmienia, jest kompresja treści omówiona we~wniosku drugim.

\item Przewaga protokołów binarnych skaluje się z~opóźnieniem sieci, a~nie z~szybkością serwera, ponieważ wyraża się liczbą zaoszczędzonych obiegów pakietu. W~komunikacji przez łącze rozległe o~opóźnieniu 26~ms jest znacząca, natomiast w~sieci lokalnej stałaby się pomijalna.

\item W~operacjach zapisu nie wykryto stałej przewagi jednego wariantu. Dla dziesięciu i~pięćdziesięciu pozycji zamówienia mediany wynosiły 83,5--165,8~ms. Przedziały ufności dla różnic REST--Direct obejmowały zero, co~nie stanowi dowodu równoważności protokołów. Interpretację ograniczała zmienność związana z~obsługą transakcji bazodanowej.

\item W~scenariuszach Echo i~odczytu z~ciepłą pamięcią podręczną wariant Direct skracał medianę czasu o~około 0,4--9,2\% względem Envoy. Nie wszystkie różnice były rozstrzygnięte statystycznie, a~przy zimnej pamięci podręcznej kolejność wariantów nie była stała.

\item Zmienność pomiarów okazała się silnie zależna od~udziału warstwy danych w~czasie odpowiedzi: od~poniżej jednego procenta w~scenariuszach izolujących narzut protokołu do~ponad pięćdziesięciu procent w~scenariuszach wymuszających zapytanie do~bazy oraz w~operacjach zapisu. Wprowadzenie pamięci podręcznej obniża tę~zmienność.

\item Przepustowość stanowiska pomiarowego przestawała rosnąć pomiędzy dwudziestoma a~trzydziestoma równoległymi klientami, podczas gdy opóźnienie nadal wzrastało. Wspólny pułap trzech wariantów wskazywał na~ograniczenie po~stronie klienta. Z~tego powodu główna macierz obejmowała najwyżej dwudziestu klientów.
\end{enumerate}

Postawiony cel pracy uznano za~osiągnięty: zaprojektowano i~zrealizowano działający system mikroserwisowy, przeprowadzono na~nim zaplanowany zestaw badań oraz udzielono odpowiedzi na~pytanie o~relację wydajnościową badanych protokołów. Odpowiedź ta~okazała się warunkowa, a~sformułowany warunek, czyli rozmiar odpowiedzi w~odniesieniu do~okna odbiorczego protokołu TCP, jest wielkością mierzalną i~możliwą do~sprawdzenia w~konkretnym zastosowaniu. Oceniono wszystkie pięć hipotez. Hipoteza H1 wymagała doprecyzowania charakteru zależności od~rozmiaru odpowiedzi. Wyniki nie pozwoliły potwierdzić równoważności wariantów w~H4, a~H5 uzyskała częściowe wsparcie.

Wyniki wspierają tezę postawioną we~wstępie. Korzyść gRPC-Web zależała od~rozmiaru odpowiedzi i~pojawiała się dla dużych, nieskompresowanych ładunków. Przy małych odpowiedziach różnice były niewielkie, a~w~operacjach zapisu nie uzyskano stałej przewagi żadnego wariantu. Analiza objętościowa wskazuje, że~kompresja może dodatkowo ograniczyć różnicę, co wymaga potwierdzenia pomiarem czasu odpowiedzi.

Zaprojektowany system ma~następujące cechy:

\begin{itemize}
\item równoległa obsługa trzech protokołów komunikacji w~ramach jednej aplikacji klienckiej, umożliwiająca ich porównanie przy identycznej logice biznesowej, tych~samych danych oraz niezmienionej ścieżce prezentacji,
\item podejście Contract-First, w~którym pliki \texttt{.proto} stanowią jedyne źródło prawdy, a~kod serwera (C\#) i~klienta (TypeScript) jest z~nich generowany,
\item kompleksowa obserwowalność obejmująca metryki, ślady rozproszone oraz pulpity wizualizacyjne, umożliwiająca korelację wyników eksperymentu z~wewnętrznym stanem systemu.
\end{itemize}

Kierunki dalszego rozwoju pracy obejmują:

\begin{itemize}
\item rozdzielenie stanowiska pomiarowego na~wiele maszyn klienckich, co~pozwoliłoby przenieść granicę wiarygodności opisaną w~ostatnim z~powyższych wniosków i~objąć pomiarem poziomy współbieżności rzędu setek równoległych klientów,
\item powtórzenie eksperymentu przy różnych wartościach opóźnienia łącza, co~umożliwiłoby empiryczną weryfikację przewidywanej proporcjonalności przewagi protokołów binarnych do~opóźnienia sieci,
\item powtórzenie macierzy z~włączoną kompresją ciała odpowiedzi na~wszystkich ścieżkach, co~pozwoliłoby zweryfikować pomiarowo przedstawione oszacowanie objętościowe i~ustalić, czy w~takiej konfiguracji różnice pomiędzy protokołami pozostają wykrywalne,
\item wdrożenie systemu na~platformie Kubernetes (AKS) z~automatycznym skalowaniem horyzontalnym i~badanie wpływu liczby replik na~wydajność.
\end{itemize}
